\notepage{22}{Uncertainty and Measurements as Filters}
For a normalized state $\ket{\psi}$, write
$\langle A\rangle=\mel{\psi}{A}{\psi}$. The uncertainty of an observable is
\[
 \Delta\Omega=\sqrt{\bigl\langle(\Omega-\langle\Omega\rangle)^2\bigr\rangle}.
\]
Expanding the square gives
\begin{align*}
 \bigl\langle(\Omega-\langle\Omega\rangle)^2\bigr\rangle
 &=\langle\Omega^2-2\langle\Omega\rangle\Omega+\langle\Omega\rangle^2I\rangle\\
 &=\langle\Omega^2\rangle-2\langle\Omega\rangle^2+\langle\Omega\rangle^2\\
 &=\langle\Omega^2\rangle-\langle\Omega\rangle^2,
\end{align*}
so $\Delta\Omega=\sqrt{\langle\Omega^2\rangle-\langle\Omega\rangle^2}$.
\smallheading{A single filter}
\begin{center}
\begin{tikzpicture}[>=Stealth,baseline]
 \node (a) at (0,0) {$\ket{\psi}$};
 \node[draw,minimum width=0.7cm,minimum height=0.7cm] (b) at (2,0) {$\Omega$};
 \node (c) at (4,0) {$\ket{\omega}$};
 \draw[->] (a)--(b);\draw[->] (b)--(c);
 \node at (2,-0.8) {$P(\omega)=\mel{\psi}{P_\omega}{\psi}$};
\end{tikzpicture}
\end{center}
The output $\ket{\omega}$ denotes the normalized selected state; the
simple ket notation assumes a nondegenerate outcome.
\smallheading{Two successive filters}
\begin{center}
\begin{tikzpicture}[>=Stealth]
 \node (a) at (0,0) {$\ket{\psi}$};
 \node[draw] (b) at (1.5,0) {$\Omega$};
 \node (c) at (3,0) {$\ket{\omega}$};
 \node[draw] (d) at (4.5,0) {$\Lambda$};
 \node (e) at (6,0) {$\ket{\lambda}$};
 \draw[->] (a)--(b);\draw[->] (b)--(c);\draw[->] (c)--(d);\draw[->] (d)--(e);
\end{tikzpicture}
\end{center}
The first probability is $P(\omega)=\mel{\psi}{P_\omega}{\psi}$;
the conditional probability for the second result is
$P(\lambda\mid\omega)=\mel{\omega}{P_\lambda}{\omega}$.
When can the second filter leave the first eigenstate unchanged, as in
\[
 \ket\psi\ \xrightarrow{\ \Omega\ }\ \ket\omega
 \ \xrightarrow{\ \Lambda\ }\ \ket\omega\ ?
\]
For simultaneous eigenstates,
\[
 \Omega\ket{\omega,\lambda}=\omega\ket{\omega,\lambda},\qquad
 \Lambda\ket{\omega,\lambda}=\lambda\ket{\omega,\lambda},
 \qquad \comm{\Omega}{\Lambda}\ket{\omega,\lambda}=0.
\]
If a complete common eigenbasis exists, $\comm{\Omega}{\Lambda}=0$:
the observables are compatible. Noncommuting observables are incompatible.
\editorial{Commutation is the finite-dimensional criterion used here.
For unbounded operators, domains and the relevant spectral commutation
conditions must also be specified.}
